%!TEX root = ../main.tex
% =============================================================
%  CAPÍTULO 3 — MARCO TEÓRICO Y ESTADO DEL ARTE
%                                       (Entrega 3 · 10 % · 28/7)
% =============================================================

\chapter{Síntesis de la literatura consultada}\label{cap:marco-teorico}

\begin{guia}[Guía de la cátedra: qué se evalúa en este capítulo]
\begin{enumerate}
    \item \textbf{Estado del arte}: investigaciones nacionales,
          internacionales y productos similares, por separado. De cada
          trabajo: qué hizo y qué resultados obtuvo, con métricas.
    \item \textbf{Brechas identificadas}: qué falta en lo existente y
          justifica \emph{este} proyecto. Es lo que separa un estado del
          arte de un listado. Incluir también los resultados negativos de
          la literatura.
    \item \textbf{Bases teóricas}: desarrollan las variables del título y su
          relación. Cada concepto termina explicando cómo se usa en el
          proyecto (teoría instrumental, no de relleno). Citas canónicas
          antes que tutoriales.
    \item \textbf{Marco conceptual}: conceptos que se nombran pero no se
          desarrollan. Los términos del glosario se marcan con
          \verb|\gls{}|, por ejemplo \gls{g:mvp} y \gls{prototipo}.
    \item \textbf{Calidad bibliográfica}: papers, tesis y libros de
          editoriales serias, con datos completos. Wikipedia y blogs restan.
    \item \textbf{Entrevistas a especialistas} (si las hay): fecha, nombre,
          cargo y aporte; transcripción en el apéndice~\ref{apx:entrevista}.
\end{enumerate}
Errores que la cátedra marca: falta de coherencia interna, bibliografía
superficial, información irrelevante o redundante. Es uno de los
capítulos más extensos (orientativo: 20 a 30 páginas).
\end{guia}

\section{Estado del arte}\label{sec:estado-del-arte}

\subsection{Investigaciones nacionales}

\citet{torino2023}, en su tesis de maestría de la Universidad Torcuato
Di Tella, desarrolló un modelo predictivo de deserción escolar para
Argentina utilizando datos socioeconómicos y habitacionales de la
Encuesta Permanente de Hogares, alcanzando un desempeño del 87,4\,\%
bajo la métrica AUC-ROC. Un hallazgo relevante para este proyecto es que
la autora señala la ausencia, a nivel nacional, de un sistema de
seguimiento educativo que permita identificar tempranamente a los
estudiantes en riesgo para brindarles un acompañamiento personalizado
--la misma carencia estructural de personalización que \sistema{} busca
atender desde el lado de la oferta de tutorías.

Un estudio sobre ruralidad y educación en Argentina \citep{olea}
cuantifica la brecha urbano-rural en la finalización del nivel
secundario: mientras que en Tierra del Fuego la diferencia es inferior a
10 puntos porcentuales, en provincias del norte del país --Chaco,
Corrientes, Misiones, Santiago del Estero, Jujuy y Tucumán-- la brecha
supera los 30 puntos porcentuales entre estudiantes urbanos y rurales de
15 a 17 años. El mismo trabajo documenta que algunas jurisdicciones ya
recurrieron históricamente a esquemas de ``escuelas mediadas por
tecnología con docentes en las ciudades y estudiantes en sus localidades
acompañados por tutores'' como estrategia para ampliar el acceso --un
antecedente conceptual directo del modelo remoto que propone \sistema{}.

Ambos antecedentes empíricos evidencian, por un lado, la necesidad
crítica de implementar sistemas de seguimiento personalizado y, por el
otro, validan estructuralmente la eficacia de los modelos mediados por
tecnología para sortear las barreras de dispersión geográfica en el
territorio nacional.

\subsection{Investigaciones internacionales}

Respecto del impacto de la tutoría personalizada, una revisión
sistemática con meta-análisis de tres niveles \citep{sciencedirect2022}
sintetizó evidencia experimental sobre tutoría privada y encontró
tamaños de efecto de entre 0,42 y 0,67 sobre el rendimiento académico,
con la región y la materia como las variables moderadoras más
influyentes entre los trece factores analizados. En la misma dirección,
un trabajo reciente sobre tutoría bajo demanda \citep{ondemand2026}
sostiene que la tutoría es una de las intervenciones educativas más
efectivas para mejorar los resultados de aprendizaje, superando en
muchos casos a otras intervenciones como la reducción del tamaño de
clase. Ambos antecedentes respaldan empíricamente la Variable~2 del
marco teórico de este proyecto (sección~\ref{sec:variable2}).

En el plano tecnológico, un trabajo de la Universidad de Ottawa
\citep{isotropy2026} propuso un sistema de recomendación semántica de
cursos universitarios basado en aprendizaje contrastivo sobre
representaciones BERT, aplicado a un catálogo de más de 500 cursos de
ingeniería. El antecedente es relevante porque identifica y corrige un
problema técnico conocido de los \gls{embeddings} tradicionales --la
anisotropía, que hace que las representaciones de distintos cursos
resulten artificialmente similares entre sí-- y constituye el precedente
más cercano encontrado al motor de \gls{matching} que \sistema{} plantea
implementar sobre perfiles de alumnos y tutores. En la misma línea, un
estudio publicado en el \emph{Journal of Applied Informatics and
Computing} \citep{paperrec2025} implementó un sistema de recomendación de
artículos científicos con sentence-transformers y similitud de coseno,
validado con una prueba de usuario de 45 personas en la que el 95,5\,\%
consideró relevantes las recomendaciones recibidas --un antecedente
metodológico útil para diseñar la validación del propio motor de
\sistema{}. Un tercer trabajo \citep{beefore2024} advierte que, hasta el
momento, el uso de sentence-transformers en sistemas de recomendación se
limita mayormente a generar información auxiliar para resolver el
problema de arranque en frío, y no a operar como mecanismo central de
recomendación --lo cual señala un espacio de aporte relativamente
inexplorado que coincide con el uso que \sistema{} le daría a esta
tecnología.

Sobre los sistemas de reputación y confianza en plataformas bilaterales,
un capítulo del libro \emph{Reengineering the Sharing Economy}
\citep{reputation} desarrolla la teoría detrás de los mecanismos de
reputación en mercados online y plantea como pregunta central de diseño
si la reputación debe ser bidireccional --como en Airbnb-- o
unidireccional --como en eBay o Amazon--, señalando que ambos
marketplaces resuelven de forma distinta el problema de la información
asimétrica entre las partes. Esta discusión respalda directamente la
decisión de diseño de \sistema{} de implementar un sistema de
calificación bidireccional entre alumnos y tutores.

Finalmente, en cuanto a la viabilidad técnica de la tutoría virtual, un
trabajo presentado en una conferencia de la ACM \citep{webrtc2025} sobre
un sistema de enseñanza remota basado en \gls{webrtc} para cursos de
inglés reportó, tras la optimización de la transmisión de audio y video
y la implementación de mecanismos de redundancia multicanal, una tasa de
pérdida de video de apenas 0,15\,\% y una latencia de pizarra de
8,2~milisegundos en un entorno de red educativa --cifras que evidencian
que es técnicamente viable sostener sesiones de videollamada estables
incluso en condiciones de conectividad variable, tal como requiere la
hipótesis central del proyecto.

\subsection{Productos y plataformas similares}

La tabla~\ref{tab:comparativa} sintetiza el relevamiento comercial de
las cinco plataformas más representativas del sector.

\begin{landscape}
    \small
    \begin{xltabular}{\linewidth}{l L L L L L}
        \caption{Comparación de plataformas de tutorías relevadas}
        \label{tab:comparativa} \\
        \toprule
        \tb{Plataforma} & \tb{Modelo de matching} & \tb{Aula virtual integrada} & \tb{Pago dentro de la plataforma} & \tb{Verificación de tutores} & \tb{Alcance / foco geográfico} \\
        \midrule
        \endfirsthead
        \toprule
        \tb{Plataforma} & \tb{Modelo de matching} & \tb{Aula virtual integrada} & \tb{Pago dentro de la plataforma} & \tb{Verificación de tutores} & \tb{Alcance / foco geográfico} \\
        \midrule
        \endhead
        \bottomrule
        \endfoot
        Preply \citep{preply} & Filtros manuales (materia, precio, disponibilidad); sin \gls{matching} semántico & Sí, parcial (Preply Classroom, hoy limitada mayormente a tutores de inglés) & Sí, por suscripción -- comisión de 100\,\% en la primera clase y de 18\,\% a 33\,\% en las siguientes & Reseñas de alumnos + video de presentación del tutor & Global; sin foco en zonas de baja densidad de oferta \\
        Wyzant \citep{wyzant} & Búsqueda por filtros & No nativa & Sí & Reseñas & Estados Unidos; más de 80.000 tutores y 4 millones de alumnos \\
        Superprof \citep{superprof} & Ninguno -- el alumno contacta manualmente a cada profesor & No & No -- el pago se coordina fuera de la plataforma & Reseñas de alumnos & Global, incluida Argentina; sin comisión para el profesor \\
        Tusclases \citep{tusclases} & Filtros (materia, precio, ubicación); sin \gls{matching} automático & No & No & Reseñas + título declarado por el profesor & España y Latinoamérica; portal líder en Argentina \\
        profe.com \citep{profecom} & Asignación manual por parte de un equipo curado, no es un \gls{marketplace} abierto & Sí, con pizarra digital & Sí & Selección propia del equipo docente de la plataforma & España; sin presencia en Argentina \\
    \end{xltabular}
\end{landscape}

Un dato relevante encontrado durante este relevamiento es que
\citep{profecom} ya ofrece a las familias un panel de seguimiento
con ``informes personalizados generados por IA con feedback detallado''
sobre el progreso del alumno --el antecedente comercial más cercano
encontrado al asistente conceptual post-sesión que \sistema{} plantea en
su alcance (sección~\ref{sec:alcance}), aunque aplicado sobre un modelo
de plantilla docente curada y no sobre un \gls{marketplace} abierto de
tutores independientes.

\textbf{Síntesis.} Ninguna de las plataformas relevadas combina
simultáneamente las cuatro piezas que propone \sistema{}: (1) un motor de
\gls{matching} basado en similitud semántica en lugar de filtros
manuales, (2) un aula virtual nativa disponible para todas las materias
y no solo para determinados perfiles de tutor, (3) un modelo de pago
protegido dentro de la plataforma, y (4) un foco explícito en ampliar la
oferta hacia zonas geográficas donde hoy es escasa o inexistente. Los
competidores de mayor escala en Argentina (Tusclases, Superprof)
replican en esencia el modelo de directorio de contactos que el
capítulo~\ref{cap:introduccion} describe como parte del problema
--búsqueda por filtros, contacto manual, sin protección de pago--
mientras que los competidores con aula virtual y pagos integrados
(Preply) no tienen foco geográfico en el país ni resuelven el
\gls{matching} de forma semántica. Ese espacio vacío es, en definitiva,
el que este proyecto busca ocupar.

\subsection{Oportunidades de diferenciación adicionales}

A partir del relevamiento de la sección anterior, se identificaron
cuatro funcionalidades que ninguna de las plataformas competidoras
ofrece y que se proponen incorporar al alcance del proyecto: el sistema
de detección de contenido inapropiado en tiempo real, la reputación con
señales implícitas, el precio dinámico por poder adquisitivo regional y
el modo de baja conectividad.

\begin{description}
    \item[Sistema de detección de contenido inapropiado en tiempo real
    (``kill-switch'' de seguridad).] Ninguna de las plataformas relevadas
    implementa moderación automática del contenido de la videollamada en
    sí misma. Existe, sin embargo, un antecedente académico directo:
    \citet{liang2013} desarrollaron SafeVchat, un sistema de detección
    automática de contenido obsceno para servicios de videollamada
    aleatoria, que combina múltiples detectores visuales mediante teoría
    de Dempster-Shafer para decidir en tiempo real si un usuario debe ser
    bloqueado. Reportaron que la combinación del sistema automático con
    moderación humana redujo el contenido ofensivo del 33,08\,\% al
    3,49\,\%. Dado que una porción significativa de los usuarios de
    \sistema{} serán menores de edad, se propone un mecanismo que analice
    localmente (\emph{edge computing}) los fotogramas mediante un modelo
    de clasificación y que ante una detección positiva corte
    automáticamente la transmisión y bloquee la cuenta preventivamente.
    El mecanismo técnico completo, incluida la verificación de edad del
    tutor mediante \gls{ocr} y DNI, se desarrolla en la
    sección~\ref{sec:seguridad}.

    \textbf{Salvedades a considerar:} ningún clasificador automático
    tiene 0\,\% de falsos positivos/negativos; es una capa adicional, no
    una garantía absoluta. El procesamiento en el dispositivo (por
    ejemplo, TensorFlow.js) evita enviar video a un servidor, facilitando
    el cumplimiento de la Ley de Protección de Datos Personales
    (Ley~25.326). El proceso de reactivación manual debe contemplar la
    notificación a un adulto responsable. Al ser un desarrollo complejo
    (visión por computadora en tiempo real), debe evaluarse su viabilidad
    dentro de los plazos del \gls{mvp}.

    \item[Reputación con fricción reducida (señales implícitas).] La
    literatura \citep{reputation} señala que las calificaciones
    explícitas (estrellas) son propensas a sesgos. Se propone combinar
    una acción mínima (reacción rápida) con señales de comportamiento
    capturadas automáticamente: tasa de sesiones completadas,
    puntualidad, y la tasa de recontratación del mismo tutor por parte
    del mismo alumno (preferencia revelada).

    \item[Precio dinámico por poder adquisitivo regional.] Se propone que
    la plataforma sugiera un rango de precio de referencia ajustado al
    poder adquisitivo de la provincia o localidad del alumno, calculado a
    partir de los datos transaccionales de la plataforma, manteniendo la
    decisión final en el tutor.

    \item[Modo de baja conectividad.] Se propone una degradación
    progresiva de calidad (video $\rightarrow$ solo audio
    $\rightarrow$ mensajería de texto) ante caídas de ancho de banda,
    sumado a la retención temporal de un buffer de video rotativo de 30
    segundos, el cual se persiste y envía como evidencia únicamente si se
    dispara el sistema de seguridad (kill-switch).
\end{description}

\subsection{Marco legal}

El diseño de la relación contractual entre la plataforma, el adulto
responsable y el menor se apoya en el sistema de capacidad progresiva
del Código Civil y Comercial de la Nación, artículos 25 y 26, que
establece que los adolescentes de entre 13 y 18 años tienen capacidad de
ejercicio creciente, pero que los contratos que los comprometan
patrimonialmente requieren la representación o el consentimiento de los
progenitores según los artículos 682 y 690, salvo los de escasa cuantía
de la vida cotidiana previstos en el artículo 684 bis, categoría en la
que un paquete de clases particulares no encuadra. De ahí surge la
decisión de que la cuenta contratante y pagadora sea siempre la del
adulto responsable.

El mismo criterio se refuerza con un antecedente impositivo: el modelo
de retención que aplican plataformas comparables en Argentina, como
Uber, Rappi, PedidosYa y Airbnb, donde AFIP obliga al agente de pago a
retener hasta un 28\,\% de lo abonado a quien no esté inscripto en el
régimen de Monotributo, funciona como antecedente directo de la
estrategia de formalización tributaria de los tutores de \sistema{},
desarrollada en el capítulo~\ref{cap:economica}. A esto se suma la
jurisprudencia laboral argentina sobre la relación de dependencia
encubierta de repartidores de plataformas de delivery, que constituye un
antecedente de riesgo legal relevante para el vínculo entre \sistema{} y
sus tutores, documentado como riesgo en el capítulo~\ref{cap:riesgos}.

\section{Bases teóricas}\label{sec:bases-teoricas}

\subsection{Variable 1: matching semántico e inteligencia artificial
aplicada a sistemas de recomendación}\label{sec:variable1}

El motor de \gls{matching} de \sistema{} se apoya conceptualmente en dos
cuerpos teóricos: los sistemas de recomendación (\emph{recommender
systems}), que buscan predecir la afinidad entre dos entidades a partir
de sus atributos; y el procesamiento de lenguaje natural (\gls{nlp})
aplicado a la representación semántica de texto mediante
\gls{embeddings}. Un \gls{embeddings} es una representación vectorial de
un texto en un espacio de alta dimensionalidad. Modelos como los
sentence-transformers generan vectores a partir de oraciones completas,
superando las búsquedas por palabras clave. Sobre este espacio,
estructuras como FAISS (Facebook AI Similarity Search) permiten
búsquedas de similitud eficientes.

Como se detalló en la sección~\ref{sec:estado-del-arte}, el trabajo de
la Universidad de Ottawa \citep{isotropy2026} validó esta técnica en la
recomendación de cursos, corrigiendo el problema de anisotropía. Este
antecedente es directamente aplicable a \sistema{}, donde los perfiles
comparten un vocabulario acotado (materias, niveles). La literatura
\citep{beefore2024} identifica que el uso de sentence-transformers como
mecanismo central sigue siendo poco explorado, confirmando que la
implementación de \sistema{} representa un aporte que excede la
aplicación estándar.

\subsection{Variable 2: acceso a educación personalizada y brecha de
aprendizaje}\label{sec:variable2}

La base teórica proviene de la investigación sobre enseñanza
individualizada de Carroll \citep{carroll1963} y Bloom
\citep{bloom1968, bloom1984}. Carroll planteó que el aprendizaje depende
de la relación entre el tiempo necesario para dominar un contenido y el
tiempo efectivamente destinado. Bloom desarrolló la teoría del
\emph{mastery learning}, sosteniendo que la mayoría puede alcanzar
niveles altos de dominio con instrucción personalizada.

El aporte central es el ``problema de los 2 sigma'' \citep{bloom1984}:
alumnos con tutoría individual rinden, en promedio, dos desviaciones
estándar por encima de alumnos en instrucción grupal convencional. La
propuesta de \sistema{} responde tecnológicamente a este problema: busca
hacer que la tutoría individual sea accesible a escala, superando la
barrera de disponibilidad geográfica. Esto es consistente con el
meta-análisis de tutoría privada \citep{sciencedirect2022} y el trabajo
sobre tutoría bajo demanda \citep{ondemand2026}, que la ubican entre las
intervenciones más efectivas.

\subsection{Relación entre las variables}

El vínculo se explica a través de la teoría económica de los mercados
bilaterales (\emph{two-sided markets}) \citep{rochet2006}. Un mercado
bilateral conecta a dos grupos (alumnos y tutores) donde la participación
de uno aumenta el valor para el otro. El mayor riesgo es el problema del
\gls{coldstart}, señalado en el capítulo~\ref{cap:introduccion}.

En un mercado geográficamente disperso (Variable~2), el problema del
\gls{coldstart} se agrava. La función del \gls{matching} semántico
(Variable~1) es el mecanismo que permite que esa masa crítica, aunque
dispersa, se encuentre eficientemente a escala nacional. Al no depender
de coincidencias exactas ni proximidad, el motor permite a \sistema{}
operar como un único mercado denso. Esta es la hipótesis técnica que
sostiene el proyecto.

\section{Marco conceptual}

\sistema{} se ubica dentro del sector \gls{edtech} y, específicamente, en
los \gls{marketplace} educativos (dimensionado en el
capítulo~\ref{cap:introduccion}).

En el plano de la infraestructura, el proyecto se apoya en \gls{webrtc},
un estándar abierto para comunicación en tiempo real en navegadores.
Sobre este, \sistema{} utiliza LiveKit, un framework open-source que
simplifica la implementación de salas escalables. Para sostener el estado
en tiempo real de esas salas, se utiliza Redis (base de datos en memoria
clave-valor), necesaria por sus exigencias de baja latencia.

En el circuito de pagos, se apoya en un modelo de \gls{escrow}
implementado a través de MercadoPago Marketplace, reteniendo los fondos
hasta la realización de la clase, reduciendo el riesgo de fraude.

Sobre la confianza, \sistema{} adopta un sistema de reputación
bidireccional, respaldado por la literatura \citep{reputation}. Para
medir la satisfacción general, se contempla el uso del \gls{nps}.
Finalmente, el proyecto se desarrolla bajo un esquema de Scrum adaptado
(capítulo~\ref{cap:metodologia}).

\section{Entrevistas a especialistas}\label{sec:entrevistas}

Se proponen dos perfiles de especialista complementarios cuyos guiones
completos se transcriben en el apéndice~\ref{apx:entrevista}:
\textbf{Perfil 1}, docente o gestor/a institucional, para indagar sobre
el rol del apoyo escolar en la prevención de la repitencia, las
diferencias urbanas-rurales en el acceso a tutorías y los recaudos
pedagógicos y de seguridad de una tutoría virtual; y \textbf{Perfil 2},
psicopedagogía / vínculo y seguridad con menores, para abordar el
impacto vincular del acompañamiento personalizado, las señales de alerta
en entornos remotos y los errores a evitar en el diseño del sistema de
reputación.

\textbf{[PENDIENTE:} una vez realizada la entrevista, se agregará acá la
síntesis de los resultados con fecha, nombre, cargo y aporte de cada
persona entrevistada. La evidencia detallada se conserva en el
apéndice~\ref{apx:entrevista}\textbf{]}